\documentclass[../../../include/open-logic-section]{subfiles}

\begin{document}

\olfileid[pt]{sth}{ordinals}{opps}
\olsection{Ordinais Sucessores e Limites}

Nos próximos capítulos, usaremos muito os ordinais. Por isso convém
introduzir algumas noções simples.

\begin{defn}
Para qualquer ordinal $\alpha$, o seu \emph{sucessor} é $\ordsucc{\alpha}
=\alpha \cup \{\alpha\}$. Dizemos que $\alpha$ é um ordinal
\emph{sucessor} se $\ordsucc{\beta} = \alpha$ para algum ordinal~$\beta$.
Dizemos que $\alpha$ é um ordinal \emph{limite} sse $\alpha$ não é nem
vazio nem um ordinal sucessor.
\end{defn}
\noindent
O resultado seguinte mostra que esta é a noção certa de
\emph{sucessor}:

\begin{prop}
Para qualquer ordinal $\alpha$:
\begin{enumerate}
	\item $\alpha \in \ordsucc{\alpha}$;
	\item $\ordsucc{\alpha}$ é um ordinal;
	\item não existe ordinal $\beta$ tal que $\alpha \in \beta \in
	\ordsucc{\alpha}$.
	\end{enumerate}
\end{prop}

\begin{proof}
Trivialmente, $\alpha \in \alpha \cup \{\alpha\} = \ordsucc{\alpha}$.
Igualmente, $\ordsucc{\alpha}$ é um conjunto transitivo de ordinais, e
portanto um ordinal por \olref[basic]{corordtransitiveord}.
E é impossível que $\alpha \in \beta \in \ordsucc{\alpha}$,
pois então ou $\beta \in \alpha$ ou $\beta = \alpha$,
contradizendo \olref[basic]{ordordered}.
\end{proof}
\noindent
Isto autoriza também uma variante de prova por indução transfinita:

\begin{thm}[Indução Transfinita Simples]\ollabel{simpletransrecursion}
Seja $\phi(x)$ uma fórmula tal que:
\begin{enumerate}
	\item $\phi(\emptyset)$; e
	\item para qualquer ordinal $\alpha$, se $\phi(\alpha)$ então
	$\phi(\ordsucc{\alpha})$; e 
	\item se $\alpha$ é um ordinal limite e $(\forall \beta \in
	\alpha)\phi(\beta)$, então $\phi(\alpha)$.
\end{enumerate}
Então $\forall \alpha \phi(\alpha)$.
\end{thm}

\begin{proof}
Provamos a contrapositiva. Logo, suponha que existe algum ordinal para o qual
vale $\lnot\phi$; seja $\gamma$ o menor tal ordinal. Então ou
$\gamma = \emptyset$, ou $\gamma = \ordsucc{\alpha}$ para algum $\alpha$
tal que $\phi(\alpha)$; ou $\gamma$ é um ordinal limite e $(\forall
\beta \in \gamma)\phi(\beta)$.
\end{proof}
\noindent
Um último bocado de notação será útil mais adiante:

\begin{defn}\ollabel{defsupstrict}
Se $X$ é um conjunto de ordinais, então $\supstrict(X) = \bigcup_{\alpha
\in X} \ordsucc{\alpha}$.
\end{defn}
\noindent
Aqui, ``lsub'' abrevia ``least strict upper bound'' (majorante superior estrito mínimo).\footnote{Alguns livros usam ``$\text{sup}(X)$'' para isto. Mas outros usam
``$\text{sup}(X)$'' para o majorante superior \emph{não estrito} mínimo, isto é,
simplesmente $\bigcup X$. Se $X$ tem um maior elemento, $\alpha$, estas
noções divergem: o majorante superior estrito mínimo é
$\ordsucc{\alpha}$, enquanto o majorante superior não estrito mínimo é
apenas $\alpha$.} O resultado seguinte explica isto:

\begin{prop}
Se $X$ é um conjunto de ordinais, $\supstrict(X)$ é o menor ordinal
estritamente maior do que todo o ordinal em $X$.
\end{prop}

\begin{proof}
Seja $Y = \Setabs{\ordsucc{\alpha}}{\alpha \in X}$, de modo que
$\supstrict(X) = \bigcup Y$. Como os ordinais são transitivos e todo o
membro de um ordinal é um ordinal, $\supstrict(X)$ é um conjunto transitivo
de ordinais, e portanto um ordinal por
\olref[basic]{corordtransitiveord}. 

Se $\alpha \in X$, então $\ordsucc{\alpha} \in Y$, logo $\ordsucc{\alpha}
\subseteq \bigcup Y = \supstrict(X)$, e portanto $\alpha \in
\supstrict(X)$. Logo $\supstrict(X)$ é estritamente maior do que todo o
ordinal em $X$.

Inversamente, se $\alpha \in \supstrict(X)$, então $\alpha \in
\ordsucc{\beta} \in Y$ para algum $\beta \in X$, de modo que $\alpha \leq
\beta \in X$. Logo $\supstrict(X)$ é o majorante superior estrito
\emph{mínimo} de $X$.
\end{proof}

\end{document}
